\chapter{Resultados e Discussão}
\label{cap:resultados}

\section{Considerações Iniciais}

Este capítulo apresenta os resultados da comparação entre as duas versões da aplicação, de
acordo com os critérios definidos no Capítulo~\ref{cap:metodologia}: as métricas de código, que
fundamentam a análise da manutenibilidade e da complexidade de desenvolvimento, as medições de
implantação, os resultados da calibração, uma execução única do teste de estresse em cada
cenário, realizada para definir as taxas dos testes de carga, e os resultados dos testes de
escalabilidade, com cinco repetições de cada configuração.

\section{Desempenho}

A calibração consistiu em um teste de estresse de cada cenário, com a taxa aumentada em oito
degraus de um minuto e com o orçamento padrão de 4 CPUs e 4 GB para a camada de aplicação. A
Tabela~\ref{tab:calibracao} apresenta a vazão máxima sustentada por cada arquitetura, isto é, a
maior taxa em que o percentil 95 do tempo de resposta permaneceu abaixo de um segundo e a taxa
de erros abaixo de 1\%, além do componente que limitou a vazão.

\begin{table}[htbp]
\centering
\caption{Vazão máxima sustentada na calibração (resultado preliminar, execução única)}
\label{tab:calibracao}
\small
\begin{tabular}{|>{\raggedright\arraybackslash}p{1.9cm}|>{\centering\arraybackslash}p{2.0cm}|>{\raggedright\arraybackslash}p{4.0cm}|>{\centering\arraybackslash}p{2.0cm}|>{\raggedright\arraybackslash}p{4.0cm}|}
\hline
\multirow{2}{*}{\textbf{Cenário}} & \multicolumn{2}{c|}{\textbf{Monolítica}} & \multicolumn{2}{c|}{\textbf{Microsserviços}} \\ \cline{2-5}
 & \textbf{Vazão (it/s)} & \textbf{Fator limitante} & \textbf{Vazão (it/s)} & \textbf{Fator limitante} \\ \hline
Catálogo & $\geq$ 600 & Não saturou; uso máximo de 2,1 das 4 CPUs & 300 & CPU do \textit{gateway} e do serviço de produtos \\ \hline
Compra & 38 & \textit{Pool} de conexões com o banco de dados & 19 & CPU dos serviços de pedidos e de pagamentos \\ \hline
Misto & 200 & \textit{Pool} de conexões com o banco de dados & 200 & CPU do \textit{gateway} e dos serviços de produtos, pedidos e pagamentos \\ \hline
\end{tabular}
\par\vspace{4pt}{\footnotesize Fonte: Elaborada pelo autor (2026). Cada iteração do catálogo faz três requisições, e cada iteração da compra faz duas.}
\end{table}

No cenário de catálogo, composto apenas por consultas, o monólito sustentou todos os degraus
até 600 iterações por segundo, cerca de 1.800 requisições por segundo, com percentil 99 abaixo
de 4~ms, utilizando no máximo 2,1 das 4 CPUs disponíveis. A versão em microsserviços sustentou
300 iterações por segundo e, a partir de 375, o percentil 95 ultrapassou um segundo. O limite foi
imposto pelo \textit{gateway} e pelo serviço de produtos, que atingiram o teto de 0,75 CPU
atribuído a cada um, enquanto outros serviços permaneciam praticamente ociosos. Esse resultado
ilustra uma consequência da divisão do orçamento: no monólito, toda a capacidade de
processamento fica disponível para a operação que estiver em uso, enquanto nos microsserviços
cada serviço só pode usar a sua parte, e o \textit{gateway} participa de todas as requisições.

No cenário de compra, o monólito sustentou 38 iterações por segundo e a versão em
microsserviços, 19. No monólito, porém, o limite não foi o processador, do qual utilizou em
média 0,37 das 4 CPUs disponíveis, mas o \textit{pool} de conexões com o banco de dados, que tem dez conexões na
configuração padrão do Spring Boot. Como a criação do pedido ocorre dentro de uma transação que
inclui a chamada ao Mercado Pago, cada pedido mantém uma conexão ocupada durante todo o tempo de
resposta do serviço externo. Nos microsserviços, cada serviço possui o seu próprio
\textit{pool}, mas o limite foi atingido antes, pelo processador dos serviços de pedidos e de
pagamentos, cada um com 0,5 CPU.

No cenário misto, as duas arquiteturas sustentaram 200 iterações por segundo e saturaram no
degrau seguinte, mas por motivos diferentes: o monólito pelo \textit{pool} de conexões, quando
a taxa de compras chegou a 50 por segundo, e os microsserviços pelo processador dos serviços
envolvidos. Enquanto ambas se mantiveram abaixo do ponto de saturação, o percentil 95 do tempo de
resposta foi de cerca de 110~ms no monólito e de 125 a 300~ms nos microsserviços, valores
dominados pelo atraso de 100~ms do simulador de pagamento nas requisições de compra.

Após a saturação, o comportamento das arquiteturas se diferenciou. No monólito, quando o fluxo de
compra esgotou o \textit{pool} de conexões, as consultas ao catálogo também foram afetadas: o seu
percentil 95 passou de 2,2~ms para 2,9~s, porque as requisições de compra em espera ocuparam os
recursos compartilhados da aplicação. Nos microsserviços, no mesmo degrau, o percentil 95 do
catálogo ficou em 288~ms, já que os serviços de produtos e de loja não compartilham recursos com
os de pedidos e de pagamentos. Esse isolamento de falhas é uma das vantagens atribuídas aos
microsserviços por \citeonline{newman2021}.

A Tabela~\ref{tab:recursos-calibracao} apresenta o consumo médio de CPU e de memória da camada de
aplicação durante a calibração.

\begin{table}[htbp]
\centering
\caption{Consumo médio da camada de aplicação durante a calibração (resultado preliminar)}
\label{tab:recursos-calibracao}
\small
\begin{tabular}{|>{\raggedright\arraybackslash}p{2.4cm}|>{\centering\arraybackslash}p{2.8cm}|>{\centering\arraybackslash}p{2.8cm}|>{\centering\arraybackslash}p{2.8cm}|>{\centering\arraybackslash}p{2.8cm}|}
\hline
\multirow{2}{*}{\textbf{Cenário}} & \multicolumn{2}{c|}{\textbf{CPU média (núcleos)}} & \multicolumn{2}{c|}{\textbf{Memória média (MB)}} \\ \cline{2-5}
 & \textbf{Monolítica} & \textbf{Microsserviços} & \textbf{Monolítica} & \textbf{Microsserviços} \\ \hline
Catálogo & 0,99 & 1,57 & 572 & 2.098 \\ \hline
Compra & 0,37 & 1,11 & 676 & 2.389 \\ \hline
Misto & 0,61 & 2,08 & 607 & 2.484 \\ \hline
\end{tabular}
\par\vspace{4pt}{\footnotesize Fonte: Elaborada pelo autor (2026). Valores médios ao longo de todo o teste de estresse, incluindo os degraus após a saturação.}
\end{table}

Em todos os cenários, a versão em microsserviços consumiu mais recursos para atender à mesma
carga: entre 1,6 e 3,4 vezes mais CPU e de 3,5 a 4,1 vezes mais memória. A diferença de
memória decorre principalmente de cada serviço executar sua própria máquina virtual Java, com o
\textit{framework} e as bibliotecas carregados em memória sete vezes; a diferença de CPU inclui o
trabalho adicional de serialização, transmissão e validação das requisições entre os serviços.

\section{Escalabilidade}

Os testes de escalabilidade utilizaram o teste de estresse do cenário misto, com oito degraus
de trinta segundos, repetido cinco vezes para cada configuração, alternando a ordem das
arquiteturas. Em todas as configurações, considera-se como vazão máxima sustentada a taxa do
último degrau em que o percentil 95 do tempo de resposta permaneceu abaixo de um segundo e a
taxa de erros abaixo de 1\%. Como a taxa aumenta em degraus, a vazão real de saturação está entre
o valor informado e o degrau seguinte, que é 50 iterações por segundo maior nos testes com taxa
máxima de 400 e 75 maior nos testes com taxa máxima de 600.

\subsection{Escalabilidade Vertical}

A Tabela~\ref{tab:escala-vertical} apresenta os resultados com orçamentos de 2, 4 e 8 CPUs para a
camada de aplicação, mantendo a memória em 4 GB. Nos microsserviços, o orçamento é dividido
entre os serviços na mesma proporção utilizada no orçamento padrão.

\begin{table}[htbp]
\centering
\caption{Escalabilidade vertical no cenário misto}
\label{tab:escala-vertical}
\small
\begin{tabular}{|>{\raggedright\arraybackslash}p{2.6cm}|>{\centering\arraybackslash}p{1.4cm}|>{\centering\arraybackslash}p{2.3cm}|>{\centering\arraybackslash}p{2.0cm}|>{\centering\arraybackslash}p{2.0cm}|>{\centering\arraybackslash}p{2.7cm}|}
\hline
\textbf{Arquitetura} & \textbf{CPUs} & \textbf{Vazão sustentada (it/s)} & \textbf{CPU média (núcleos)} & \textbf{Memória máxima (MB)} & \textbf{Inicialização (s)} \\ \hline
\multirow{3}{*}{Monolítica} & 2 & 200 & 0,61 & 600 & 18 \\ \cline{2-6}
 & 4 & 200 & 0,61 & 681 & 17 \\ \cline{2-6}
 & 8 & 225 & 0,61 & 830 & 16 \\ \hline
\multirow{3}{*}{Microsserviços} & 2 & $<$ 50 & 1,10 & 2.705 & 107 \\ \cline{2-6}
 & 4 & 150 & 2,14 & 2.578 & 73 \\ \cline{2-6}
 & 8 & 225 & 3,08 & 2.639 & 57 \\ \hline
\end{tabular}
\par\vspace{4pt}{\footnotesize Fonte: Elaborada pelo autor (2026). A vazão sustentada foi a mesma nas cinco repetições de cada configuração. CPU média e memória máxima referem-se à camada de aplicação durante todo o teste; a inicialização é o tempo até todas as aplicações estarem prontas.}
\end{table}

O monólito praticamente não se beneficiou do aumento de processadores. Com 2 CPUs, sustentou as
mesmas 200 iterações por segundo obtidas com 4 CPUs e, com 8 CPUs, chegou a 225, um ganho de
12,5\% para o dobro de recursos, o que corresponde a uma eficiência de 0,56. O consumo médio de
CPU foi de 0,61 núcleo nos três orçamentos, e o percentil 90 das amostras ficou em torno de
0,95 núcleo, de modo que a capacidade de processamento adicional permaneceu ociosa. Esse
comportamento confirma o que foi observado na calibração: a vazão do monólito no cenário misto é
limitada pelo \textit{pool} de dez conexões com o banco de dados, ocupadas durante a chamada ao
Mercado Pago, e esse limite não depende do número de processadores.

Os microsserviços, por outro lado, foram limitados pelo processador em todos os orçamentos. Com
2 CPUs, cada serviço recebeu entre 0,25 e 0,375 CPU, e a versão não sustentou nem o primeiro
degrau, de 50 iterações por segundo: o percentil 95 das compras chegou a 12 segundos, com o
\textit{gateway} e os serviços de produtos, lojas, pedidos e pagamentos operando no limite de CPU
atribuído. Com 4 CPUs, a vazão sustentada foi de 150 iterações por segundo. Na calibração, que
utilizou degraus de um minuto, essa configuração havia sustentado 200; nas cinco repetições, o
degrau de 200 resultou em um percentil 95 médio de 1,4~s, pouco acima do critério, o que situa a
capacidade dessa configuração entre 150 e 200 iterações por segundo. Com 8 CPUs, a vazão chegou a
225, um ganho de 50\% em relação a 4 CPUs e eficiência de 0,75. Nesse orçamento, o limite passou a
ser o \textit{gateway}, que operou no seu teto de 1,5 CPU, enquanto os serviços de pedidos e de
pagamentos utilizaram cerca de metade da CPU disponível.

Com 8 CPUs, as duas arquiteturas sustentaram a mesma vazão, mas os microsserviços utilizaram, em
média, cinco vezes mais CPU e três vezes mais memória do que o monólito. A memória dos
microsserviços variou pouco entre os orçamentos, por ser determinada principalmente pelas sete
máquinas virtuais Java. O orçamento de CPU também afetou a inicialização: o tempo até todas as
aplicações estarem prontas caiu de 107 para 57 segundos nos microsserviços, enquanto no monólito
permaneceu entre 16 e 18 segundos.

\subsection{Escalabilidade Horizontal}

A Tabela~\ref{tab:escala-horizontal} apresenta os resultados com o balanceador de carga Nginx. No
monólito, a aplicação inteira foi executada em duas instâncias, cada uma com 4 CPUs e 4 GB. Nos
microsserviços, foram executadas duas instâncias do \textit{gateway} e dos serviços de produtos,
de pedidos e de pagamentos, que atingiram o limite de CPU na calibração do cenário misto, cada
instância com os mesmos limites da instância original. A eficiência é a razão entre o ganho de
vazão e o aumento do orçamento de CPU da camada de aplicação em relação à referência.

\begin{table}[htbp]
\centering
\caption{Escalabilidade horizontal no cenário misto}
\label{tab:escala-horizontal}
\small
\begin{tabular}{|>{\raggedright\arraybackslash}p{2.5cm}|>{\raggedright\arraybackslash}p{2.0cm}|>{\centering\arraybackslash}p{1.2cm}|>{\centering\arraybackslash}p{2.3cm}|>{\centering\arraybackslash}p{1.8cm}|>{\centering\arraybackslash}p{1.8cm}|>{\centering\arraybackslash}p{1.9cm}|}
\hline
\textbf{Arquitetura} & \textbf{Instâncias} & \textbf{CPUs} & \textbf{Vazão sustentada (it/s)} & \textbf{CPU média (núcleos)} & \textbf{Memória máxima (MB)} & \textbf{Eficiência} \\ \hline
\multirow{2}{*}{Monolítica} & 1 (referência) & 4 & 200 & 0,62 & 701 & -- \\ \cline{2-7}
 & 2 & 8 & 405\newline(354 a 456) & 1,32 & 1.388 & 1,01 \\ \hline
\multirow{2}{*}{Microsserviços} & 1 de cada (referência) & 4 & 150 & 2,16 & 2.586 & -- \\ \cline{2-7}
 & 2 de quatro serviços & 6,5 & 225 & 3,78 & 4.107 & 0,92 \\ \hline
\end{tabular}
\par\vspace{4pt}{\footnotesize Fonte: Elaborada pelo autor (2026). Entre parênteses, o intervalo de confiança de 95\% da média das cinco repetições; nas demais configurações, a vazão sustentada foi a mesma em todas as repetições.}
\end{table}

As execuções de referência sustentaram a mesma vazão das execuções sem balanceador, 200
iterações por segundo no monólito e 150 nos microsserviços, e o percentil 95 no primeiro degrau
variou de 109 para 110~ms no monólito e de 130 para 142~ms nos microsserviços. O custo do
balanceador foi maior nos microsserviços porque, nessa configuração, as chamadas entre os serviços
também passam por ele.

O monólito foi a arquitetura que mais se beneficiou da replicação. Com duas instâncias, a vazão
sustentada passou para uma média de 405 iterações por segundo, variando entre 375 e 450 nas cinco
repetições, o dobro da referência, com eficiência de 1,01. Diferentemente do aumento de
processadores, a replicação duplica também o \textit{pool} de conexões, já que cada instância
mantém o seu, o que remove o limite que restringia a escalabilidade vertical. Mesmo com o dobro da
carga, o uso de CPU da camada de aplicação permaneceu abaixo de um terço do orçamento, e o do banco
de dados ficou em torno de 1 das 2 CPUs disponíveis.

Nos microsserviços, a replicação de quatro serviços elevou a vazão de 150 para 225 iterações por
segundo, um ganho de 50\% para um aumento de 62,5\% no orçamento de CPU, com eficiência de 0,92.
A vazão foi a mesma obtida com 8 CPUs na escalabilidade vertical, utilizando 6,5 CPUs. Com as
réplicas, o \textit{gateway} e o serviço de pagamentos voltaram a operar no limite de CPU das duas
instâncias, e o serviço de lojas, que não foi replicado, também atingiu o seu limite de 0,5 CPU.
O aumento de memória foi de 1,5 GB, equivalente às quatro novas máquinas virtuais Java, e o tempo
até todas as aplicações estarem prontas passou de 77 para 85 segundos nos microsserviços e de 16
para 19 segundos no monólito.

\section{Manutenibilidade}

As métricas de código foram extraídas com a ferramenta CK \cite{aniche-ck} sobre o código de
produção das duas versões. A Tabela~\ref{tab:ck} apresenta os totais e as médias por classe.

\begin{table}[htbp]
\centering
\caption{Métricas de código por arquitetura (ferramenta CK)}
\label{tab:ck}
\small
\begin{tabular}{|>{\raggedright\arraybackslash}p{6.6cm}|>{\centering\arraybackslash}p{3.6cm}|>{\centering\arraybackslash}p{3.6cm}|}
\hline
\textbf{Métrica} & \textbf{Monolítica} & \textbf{Microsserviços} \\ \hline
Classes, interfaces e enumerações & 78 & 123 \\ \hline
Linhas de código (CK) & 973 & 1.470 \\ \hline
Acoplamento entre objetos (CBO) médio & 5,79 & 5,67 \\ \hline
Complexidade por classe (WMC) média & 2,73 & 2,45 \\ \hline
Resposta por classe (RFC) média & 6,17 & 5,68 \\ \hline
Falta de coesão em métodos (LCOM) média & 1,54 & 1,21 \\ \hline
Complexidade total (soma do WMC) & 213 & 301 \\ \hline
\end{tabular}
\par\vspace{4pt}{\footnotesize Fonte: Elaborada pelo autor (2026). A ferramenta não contabiliza os \textit{records} do Java, utilizados nos objetos de transferência de dados, que estão incluídos apenas nas contagens da Tabela~\ref{tab:repositorios}.}
\end{table}

As médias por classe são muito próximas nas duas versões, o que é esperado, já que as regras de
negócio e a organização interna de cada domínio são as mesmas. A diferença está no volume: a
versão em microsserviços tem 58\% mais classes, 51\% mais linhas de código e 41\% mais
complexidade total para implementar as mesmas funcionalidades. Esse acréscimo corresponde ao
código necessário para a distribuição, como os clientes HTTP, os publicadores e consumidores de
mensagens, as classes de representação de dados de outros serviços, as configurações de cada
aplicação e o \textit{API Gateway}. As médias ligeiramente menores de acoplamento e de falta de
coesão nos microsserviços refletem justamente esse grande número de classes pequenas e simples,
e não uma melhora da estrutura das regras de negócio.

\section{Complexidade de Desenvolvimento}

A Tabela~\ref{tab:repositorios} reúne as contagens extraídas dos repositórios das duas versões.

\begin{table}[htbp]
\centering
\caption{Tamanho e configuração dos repositórios}
\label{tab:repositorios}
\small
\begin{tabular}{|>{\raggedright\arraybackslash}p{7.4cm}|>{\centering\arraybackslash}p{3.2cm}|>{\centering\arraybackslash}p{3.2cm}|}
\hline
\textbf{Item} & \textbf{Monolítica} & \textbf{Microsserviços} \\ \hline
Projetos (aplicações) & 1 & 7 \\ \hline
Arquivos Java de produção & 105 & 166 \\ \hline
Linhas de código Java de produção & 2.286 & 3.483 \\ \hline
Arquivos e linhas de código de teste & 12 e 771 & 23 e 952 \\ \hline
Testes automatizados & 52 & 66 \\ \hline
Dependências declaradas nos arquivos \textit{pom.xml} & 14 & 70 \\ \hline
Arquivos de configuração e linhas (\textit{pom.xml}, \textit{application.yml}, \textit{Dockerfile}) & 4 e 173 & 28 e 1.079 \\ \hline
Arquivos e linhas de código de comunicação entre componentes & 3 e 112 & 34 e 619 \\ \hline
\end{tabular}
\par\vspace{4pt}{\footnotesize Fonte: Elaborada pelo autor (2026). As linhas não incluem linhas em branco.}
\end{table}

A versão em microsserviços tem 52\% mais linhas de código Java, cinco vezes mais dependências
declaradas e cerca de seis vezes mais linhas de configuração. O código dedicado à comunicação
entre componentes ilustra a complexidade introduzida pela distribuição: são 34 arquivos e 619
linhas, contra 3 arquivos e 112 linhas no monólito, onde essa comunicação se resume ao evento
interno de pagamento atualizado e ao seu \textit{listener}. Além disso, cada um dos sete
projetos tem suas próprias dependências, versões e configurações, o que, durante o
desenvolvimento, exigiu alinhar manualmente as versões do Spring Boot entre os serviços para que
todos utilizassem a mesma versão do monólito.

O tempo de construção, apresentado na seção seguinte, também afeta o desenvolvimento, já que
cada alteração precisa ser compilada e empacotada antes de ser testada.

\section{Implantação}

A Tabela~\ref{tab:implantacao} apresenta as medições relacionadas à implantação. Os tempos de
construção correspondem à geração das imagens Docker sem reaproveitamento de camadas, em três
repetições, e os tempos de inicialização foram medidos nas seis subidas do ambiente realizadas
durante a calibração, três para cada arquitetura, desde o comando de subida até que todas as
aplicações respondessem.

\begin{table}[htbp]
\centering
\caption{Medições de implantação}
\label{tab:implantacao}
\small
\begin{tabular}{|>{\raggedright\arraybackslash}p{7.4cm}|>{\centering\arraybackslash}p{3.2cm}|>{\centering\arraybackslash}p{3.2cm}|}
\hline
\textbf{Item} & \textbf{Monolítica} & \textbf{Microsserviços} \\ \hline
Imagens das aplicações & 1 & 7 \\ \hline
Contêineres em execução (aplicações e infraestrutura) & 3 & 11 \\ \hline
Tamanho das imagens & 593 MB & 480 a 610 MB cada; 3,9 GB no total \\ \hline
Espaço em disco, descontada a imagem base compartilhada (402 MB) & 191 MB & 1,1 GB \\ \hline
Tempo médio de construção das imagens & 52 s & 72 s \\ \hline
Tempo médio de inicialização do Spring Boot (aplicação mais lenta) & 8,8 s & 36,9 s \\ \hline
Tempo médio até todas as aplicações estarem prontas & 19 s & 76 s \\ \hline
\end{tabular}
\par\vspace{4pt}{\footnotesize Fonte: Elaborada pelo autor (2026). A contagem de contêineres não inclui o simulador do Mercado Pago, utilizado apenas nos testes.}
\end{table}

A versão em microsserviços exige sete imagens e onze contêineres, contra uma imagem e três
contêineres do monólito, já que, além das aplicações, precisa do Kafka e do Zookeeper. A
inicialização completa levou, em média, quatro vezes mais tempo. Parte dessa diferença decorre
da divisão do orçamento de CPU: cada serviço inicializa com no máximo 0,5 ou 0,75 CPU, enquanto o
monólito dispõe das 4 CPUs, e os serviços que dependem do Kafka aguardam o \textit{broker} estar
pronto.

\section{Discussão}

\section{Considerações Finais}
